\documentclass[12pt,a4paper]{article}
\usepackage[utf8]{inputenc}
\usepackage[brazil]{babel}
\usepackage[T1]{fontenc}
\usepackage{geometry}
\usepackage{graphicx}
\usepackage{hyperref}
\usepackage{listings}
\usepackage{xcolor}
\usepackage{amsmath}
\usepackage{amssymb}

\geometry{a4paper, left=3cm, right=2cm, top=3cm, bottom=2cm}

\definecolor{codegreen}{rgb}{0,0.6,0}
\definecolor{codegray}{rgb}{0.5,0.5,0.5}
\definecolor{codepurple}{rgb}{0.58,0,0.82}
\definecolor{backcolour}{rgb}{0.95,0.95,0.92}

\lstdefinestyle{mystyle}{
    backgroundcolor=\color{backcolour},
    commentstyle=\color{codegreen},
    keywordstyle=\color{magenta},
    numberstyle=\tiny\color{codegray},
    stringstyle=\color{codepurple},
    basicstyle=\ttfamily\footnotesize,
    breakatwhitespace=false,
    breaklines=true,
    captionpos=b,
    keepspaces=true,
    numbers=left,
    numbersep=5pt,
    showspaces=false,
    showstringspaces=false,
    showtabs=false,
    tabsize=2
}

\lstset{style=mystyle}

% Gramaticas e saidas nao precisam de numero de linha.
\lstdefinestyle{gram}{style=mystyle, numbers=none, basicstyle=\ttfamily\scriptsize}

% Evita linhas estouradas por trechos longos em \texttt.
\setlength{\emergencystretch}{1em}

\title{Relatório de Projeto: Compilador MiniPython\\[4pt]
\large Trabalho Prático 1: Análise Léxica e Sintática}
\author{
    Brenda Gabrielle Alves Nascimento \\
    Matrícula: 24.1.4011 \\
    BCC328 - Construção de Compiladores I - DECOM/UFOP
}
\date{\today}

\begin{document}

\maketitle

\begin{abstract}
Primeira etapa do compilador MiniPython, em Haskell: o analisador
léxico e três analisadores sintáticos, descendente recursivo
(Megaparsec), PEG (typed-peg) e LALR (Happy), que produzem a mesma AST.
\end{abstract}

% O sumario foi retirado para caber no limite de 6 paginas.
%\tableofcontents

\section{Introdução}

MiniPython é um subconjunto de Python 3 com anotações de tipo, classes
e funções anônimas, com blocos delimitados por indentação. Este
trabalho implementa a análise léxica e a sintática, acessíveis pelo
executável \texttt{mpyc}; as gramáticas estão, com mais detalhes, em
\texttt{docs/gramatica.md}.

\section{Metodologia}

\subsection{Estrutura sintática de MiniPython}

\subsubsection*{Gramática livre de contexto}

A GLC é escrita sobre os tokens do analisador léxico: os terminais em
maiúsculas são classes de token e os entre aspas são o próprio lexema;
\texttt{\{ X \}} é repetição e \texttt{[ X ]} é opcional. Os blocos são
\texttt{NEWLINE INDENT ... DEDENT}, o análogo das chaves de C.

\begin{lstlisting}[style=gram]
Programa     ::= { Decl } EOF
Decl         ::= DefClasse | Comando
DefFuncao    ::= 'def' IDENT '(' [ Params ] ')' [ '->' Tipo ] ':' Bloco
Params       ::= Param { ',' Param }
Param        ::= 'self' | IDENT [ ':' Tipo ]
DefClasse    ::= 'class' IDENT [ '(' IDENT ')' ] ':' BlocoClasse
BlocoClasse  ::= NEWLINE INDENT { Campo } { DefFuncao } [ 'pass' NEWLINE ] DEDENT
Campo        ::= IDENT ':' Tipo NEWLINE
Bloco        ::= NEWLINE INDENT Comando { Comando } DEDENT
Tipo         ::= 'int' | 'float' | 'bool' | 'str' | 'None' | IDENT [ '[' Tipo ']' ]
               | '(' [ Tipo { ',' Tipo } ] ')' '->' Tipo
Comando      ::= ComandoSimples NEWLINE | ComandoIf | ComandoWhile | ComandoFor | DefFuncao
ComandoSimples ::= 'pass' | 'break' | 'continue' | 'return' [ Expr ] | Expr
               | IDENT ':' Tipo '=' Expr | Alvo ( '=' | OpAumentado ) Expr
Alvo         ::= IDENT | ExprPosfixa '[' Expr ']' | ExprPosfixa '.' IDENT
OpAumentado  ::= '+=' | '-=' | '*=' | '/=' | '//=' | '%='
ComandoIf    ::= 'if' Expr ':' Bloco { 'elif' Expr ':' Bloco } [ 'else' ':' Bloco ]
ComandoWhile ::= 'while' Expr ':' Bloco
ComandoFor   ::= 'for' IDENT 'in' Expr ':' Bloco
Expr         ::= 'lambda' [ IDENT { ',' IDENT } ] ':' Expr | ExprOu [ 'if' ExprOu 'else' Expr ]
ExprOu       ::= ExprE { 'or' ExprE }
ExprE        ::= ExprNao { 'and' ExprNao }
ExprNao      ::= 'not' ExprNao | ExprRel
ExprRel      ::= ExprAdit [ ( '==' | '!=' | '<' | '<=' | '>' | '>=' ) ExprAdit ]
ExprAdit     ::= ExprMult { ( '+' | '-' ) ExprMult }
ExprMult     ::= ExprPot { ( '*' | '/' | '//' | '%' ) ExprPot }
ExprPot      ::= ExprUnaria [ '**' ExprPot ]
ExprUnaria   ::= '-' ExprUnaria | ExprPosfixa
ExprPosfixa  ::= Atomo { '(' [ Args ] ')' | '[' Expr ']' | '.' IDENT }
Args         ::= Expr { ',' Expr }
Atomo        ::= INT_LIT | FLOAT_LIT | STRING_LIT | BOOL_LIT | 'None' | 'self' | IDENT
               | 'int' | 'float' | 'str' | 'bool' | '[' [ Args ] ']' | '(' Expr ')'
\end{lstlisting}

A precedência e a associatividade ficam na própria forma das produções,
do nível mais fraco para o mais forte: condicional (à direita),
\texttt{or}, \texttt{and} (à esquerda), \texttt{not} (prefixo),
relacionais (não associativos), \texttt{+ -} e \texttt{* / // \%} (à
esquerda), \texttt{**} (à direita), menos unário e os pós-fixos. As
repetições \texttt{\{ op B \}} evitam a recursão à esquerda, que
impediria a implementação descendente~\cite{appel1998}, e
\texttt{a < b < c} é erro de sintaxe. Como a seção 6.3 da especificação
põe o menos unário acima de \texttt{**}, \texttt{-2 ** 2} é
\texttt{(-2) ** 2}. A lambda fica no nível de \texttt{Expr}, e não
entre os átomos: como átomo, \texttt{lambda x: x + 1} teria duas
árvores, com a soma dentro ou fora do corpo; assim o corpo vai até o
fim da expressão, e como operando a lambda precisa de parênteses. A
condição e o valor da condicional são \texttt{ExprOu}, então ela
associa à direita e \texttt{1 if 2 if 3 else 4 else 5} é erro.
\texttt{DefFuncao} também é comando, porque a regra 19 da seção 10 da
especificação prevê funções locais, e por isso não se repete em
\texttt{Decl}. O tipo lista é \texttt{IDENT [ Tipo ]}, aceito só para
\texttt{list}.

\subsubsection*{Gramática PEG}

A typed-peg~\cite{typedpeg} lê caracteres, e não tokens, então a
gramática PEG~\cite{ford2004} descreve também a parte léxica e a
indentação. Além da notação usual (\texttt{/} é a escolha ordenada,
\texttt{\&} e \texttt{!} são predicados), há as anotações de layout da
biblioteca, que seguem o modelo de Adams~\cite{adams2013}:
\texttt{e\^{}>} é um bloco mais à direita que o comando que o contém,
\texttt{|e|} exige que \texttt{e} comece alinhado com o bloco e
\texttt{e\_\textasciitilde} libera a coluna dentro de \texttt{e}. Ficam
implícitos o \texttt{Espaco} que todo token consome depois de si e o
\texttt{![a-zA-Z0-9\_]} depois de cada palavra reservada;
\texttt{Reservada} lista as 24 palavras reservadas da seção 2 do
enunciado.

\begin{lstlisting}[style=gram]
Programa     <- Brancos (Ind |Decl|)* Fim
Decl         <- DefFuncao / DefClasse / Comando
DefFuncao    <- 'def' Nome '(' (Params ')')_~ ('->' Tipo)? ':' Suite
Params       <- (Param (',' Param)*)?
Param        <- 'self' / Nome (':' Tipo)?
DefClasse    <- 'class' Nome ('(' Nome ')')? ':' FimLinha Ind Membros^>
Membros      <- |Campo| (Ind |Campo|)* (Ind |DefFuncao|)* PassFinal?
              / |DefFuncao| (Ind |DefFuncao|)* PassFinal? / |'pass'| FimLinha
PassFinal    <- Ind |'pass'| FimLinha
Campo        <- Nome ':' Tipo FimLinha
Suite        <- FimLinha Ind (|Comando| (Ind |Comando|)*)^>
Comando      <- DefFuncao / ComandoIf / ComandoWhile / ComandoFor / ComandoSimples FimLinha
ComandoIf    <- 'if' Expr ':' Suite (Ind |'elif'| Expr ':' Suite)*
                (Ind |'else'| ':' Suite)?
ComandoWhile <- 'while' Expr ':' Suite
ComandoFor   <- 'for' Nome 'in' Expr ':' Suite
ComandoSimples <- 'pass' / 'break' / 'continue' / 'return' Expr?
              / Nome ':' Tipo '=' !'=' Expr / Expr (OpAtrib Expr)?
OpAtrib      <- '=' !'=' / '+=' / '-=' / '*=' / '//=' / '/=' / '%='
Tipo         <- 'int' / 'float' / 'bool' / 'str' / 'None'
              / '(' ((Tipo (',' Tipo)*)? ')')_~ '->' Tipo / Nome ('[' (Tipo ']')_~)?
Expr         <- 'lambda' (Nome (',' Nome)*)? ':' Expr / ExprOu ('if' ExprOu 'else' Expr)?
ExprOu       <- ExprE ('or' ExprE)*
ExprE        <- ExprNao ('and' ExprNao)*
ExprNao      <- 'not' ExprNao / ExprRel
ExprRel      <- ExprAdit (('==' / '!=' / '<=' / '>=' / '<' / '>') ExprAdit)?
ExprAdit     <- ExprMult (('+' !'=' / '-' ![=>]) ExprMult)*
ExprMult     <- ExprPot (('//' !'=' / '/' !'=' / '%' !'=' / '*' ![*=]) ExprPot)*
ExprPot      <- ExprUnaria ('**' ExprPot)?
ExprUnaria   <- '-' ![=>] ExprUnaria / ExprPosfixa
ExprPosfixa  <- Atomo ('(' (Args ')')_~ / '[' (Expr ']')_~ / '.' Nome)*
Args         <- (Expr (',' Expr)*)?
Atomo        <- 'True' / 'False' / 'None' / 'self' / 'int' / 'float' / 'str' / 'bool'
              / Nome / Numero / String / '[' (Args ']')_~ / '(' (Expr ')')_~
EspacoLinha  <- ([ \t\r] / '#' [^\n]*)*
Espaco       <- ([ \t\r] / '#' [^\n]* / &('\n' .) '\n')*
FimLinha     <- (EspacoLinha ('\n' / !.) Brancos)_~
Brancos      <- (EspacoLinha '\n')*_~
Ind          <- [ \t]*_~
Fim          <- (EspacoLinha !.)_~
Nome         <- !Reservada [a-zA-Z_] [a-zA-Z0-9_]*
Reservada    <- ('and' / 'break' / ... / 'int' / 'in' / ... / 'while') ![a-zA-Z0-9_]
Numero       <- Digitos ('.' Digitos)? ([eE] [+-]? Digitos)?
Digitos      <- [0-9] ('_'? [0-9])*
String       <- '"' (Escape / [^"\\\n])* '"' / '\'' (Escape / [^'\\\n])* '\''
Escape       <- '\\' [ntr\\"']
\end{lstlisting}

A relação padrão entre tokens é \texttt{gtR}: depois do primeiro token,
o resto do comando fica à direita da coluna do comando. Isso resolve
também as quebras de linha, porque \texttt{Espaco} só consome um
\texttt{\textbackslash n} se o primeiro caractere da linha seguinte for
aceito pela relação em vigor: sob \texttt{gtR} a coluna 0 nunca é, e a
quebra encerra o comando; sob \texttt{\_\textasciitilde}, dentro de
parênteses e colchetes, ela vira espaço. Na escolha ordenada,
\texttt{'int'} vem antes de \texttt{'in'}, e os predicados nos
operadores fazem o papel do casamento mais longo do analisador léxico
(\texttt{+} não pode começar \texttt{+=}).

%\subsection{Sistema de tipos para MiniPython}
%\subsection{Inferência de tipos para MiniPython}
%\subsection{Semântica operacional para MiniPython}

\section{Arquitetura do Compilador}

O código está em \texttt{src/MiniPython/Frontend}, separado em
\texttt{Lexer}, \texttt{Parser}, \texttt{Syntax} e \texttt{Pretty}. O
executável roda sempre o analisador léxico primeiro e, conforme a opção,
um dos três parsers e a verificação de contexto da Seção~\ref{sec:ctx}.

\subsection{Análise léxica}

O scanner (\texttt{MiniPythonLexer.hs}) é escrito à mão: percorre o
texto mantendo linha e coluna, descarta comentários, expande tabulações
até a próxima coluna múltipla de 8 e emite um \texttt{NEWLINE} ao fim de
cada linha física. Uma segunda passada (\texttt{Indent.hs}) aplica o
algoritmo da seção 4.5 da especificação: compara a coluna do primeiro
token de cada linha lógica com o topo da pilha de níveis, descarta o
\texttt{NEWLINE} de linhas em branco e de linhas dentro de parênteses ou
colchetes e fecha os blocos que restam antes do \texttt{EOF}. Com as duas
passadas separadas, o scanner não guarda nenhum estado de indentação.

Os construtores do tipo dos tokens têm os nomes da tabela da seção 2.2
do enunciado (\texttt{TkDef}, \texttt{TkIdent \textquotedbl{}fib\textquotedbl}, \ldots), e cada
linha de \texttt{-{}-lex} é a posição seguida do \texttt{show} do token;
o \texttt{Show} derivado é usado só para tokens, nunca para a AST. O
\texttt{INDENT} fica no primeiro token da linha indentada, o
\texttt{DEDENT} e o erro de indentação no \texttt{NEWLINE} da linha
anterior, onde o bloco se fecha, e o \texttt{EOF} sai como
\texttt{0:0}. Os erros léxicos são caractere inválido, string não
terminada, escape inválido, nível de indentação fora da pilha e literal
numérico mal formado: \texttt{\_} fora de entre dois dígitos
(\texttt{1\_\_2}, \texttt{1\_}), ponto sem dígito (\texttt{3.}) ou
expoente sem dígitos (\texttt{1e+}).

\subsection{Análise sintática}

\paragraph{Descendente recursivo.} Escrito com Megaparsec~\cite{megaparsec}
sobre a lista de tokens, por meio de uma instância de \texttt{Stream}.
Cada não terminal da GLC é uma função, e a produção é escolhida olhando
o próximo token; só a declaração de variável e o campo de classe
precisam ver dois (\texttt{IDENT} seguido de \texttt{:}), e não há
\texttt{try}. A atribuição é lida como uma expressão que depois é
convertida em alvo, o que dispensa decidir de antemão entre
\texttt{Alvo} e \texttt{Expr}.

\paragraph{PEG.} Escrito com o quasi-quoter \texttt{pegGrammar}, com as
diretivas \texttt{\%name}, \texttt{\%stream String}, \texttt{\%param src
::\ Source} e \texttt{\%start}. A versão usada é a 0.6, porque a 0.4 só
informa que o reconhecimento falhou, sem posição, e não dá às ações o
deslocamento atual, de onde vêm as posições da AST (átomo \texttt{@}).
As regras de comando devolvem \texttt{Either Error}, porque três
restrições só podem ser verificadas na ação semântica, que não faz o
reconhecimento falhar: alvo de atribuição, parâmetro repetido e
argumento de tipo fora de \texttt{list}.

\paragraph{LALR.} A especificação do Happy~\cite{happy} não estratifica
os operadores: há uma única produção, \texttt{Operacao}, com todos
eles, e a precedência vem das diretivas, da menor para a maior:
\texttt{or}, \texttt{and}, \texttt{not}, os relacionais
(\texttt{\%nonassoc}, o que rejeita \texttt{a < b < c}), \texttt{+ -},
\texttt{* / // \%}, \texttt{**} e \texttt{NEG} (o menos unário, com
\texttt{\%prec NEG}). A lambda e a condicional ficam numa produção
\texttt{Expr} à parte, como na GLC: que a condição de uma condicional
não seja outra condicional, nenhuma diretiva expressa, porque diante do
segundo \texttt{if} em \texttt{1 if 2 if 3 ...} o único movimento
possível é deslocar. A gramática não tem conflitos.

\subsection{Regras de contexto e mensagens de erro}
\label{sec:ctx}

\texttt{return} fora de função e \texttt{break} ou \texttt{continue} fora
de laço são verificados sobre a AST (\texttt{Syntax/Validate.hs}), a
mesma para os três parsers. O parâmetro repetido é verificado por cada
parser ao fechar a lista, porque o arquivo \texttt{.err} aponta o
\texttt{)}, posição que não fica na AST; no Happy, a lista virou uma
produção própria, reduzida logo depois do \texttt{)}. Numa função
local, \texttt{return} vale, mas o laço de fora não conta: um
\texttt{break} ali é erro, como em Python.

As mensagens seguem a seção 2.1 do enunciado, com o nome do arquivo
sem o diretório; as sintáticas são \texttt{expected X, got Y}. No
descendente recursivo, X é o que a função esperava. No LALR,
\texttt{\%errorhandlertype explist} entrega os terminais aceitos no
estado do erro, e vale o mais informativo (\texttt{:}, \texttt{INDENT},
início de expressão, \texttt{NEWLINE}). No PEG, a typed-peg dá a
posição mais à direita em que alguma regra falhou, e os \texttt{\%label}
nomeiam as regras e escondem as continuações opcionais.

\subsection{Árvore de sintaxe abstrata}

A AST (\texttt{MiniPythonSyntax.hs}) tem um tipo para cada categoria:
\texttt{Program}, \texttt{Decl}, \texttt{FunDef}, \texttt{ClassDef},
\texttt{Stmt}, \texttt{LValue}, \texttt{Expr} e \texttt{Type}, e todo nó
guarda a posição do token que o inicia. Declaração com anotação
(\texttt{SVarDecl}) e atribuição (\texttt{SAssign}) são construtores
distintos, e o alvo da atribuição é um \texttt{LValue}, o que torna
impossível representar uma atribuição a uma chamada. Uma função local é
o comando \texttt{SFunDef}; as de nível superior ficam em
\texttt{DFun}, porque só elas serão generalizadas na inferência de
tipos. A impressão usa
\texttt{drawTree}: \texttt{toTree ::\ Program -> Tree String} põe no
rótulo o construtor e os campos atômicos, com os tipos escritos como no
fonte, e os blocos de \texttt{if}, \texttt{while} e \texttt{for} ficam
sob nós auxiliares em minúsculas (\texttt{then}, \texttt{elif},
\texttt{else}, \texttt{body}). Para \texttt{x: int = 1 + 2}:

\begin{lstlisting}[style=gram]
Program
|
`- SVarDecl x: int
   |
   `- EBinOp Add
      |
      +- EInt 1
      |
      `- EInt 2
\end{lstlisting}

%\subsection{Análise semântica}
%\subsection{Geração de código}

\section{Resultados e Discussão}

\subsection{Instruções de Uso}

Com GHC 9.6 e cabal: \texttt{cabal update} (a typed-peg 0.6 vem do
Hackage) e \texttt{cabal build}. Para rodar,
\texttt{cabal run mpyc -{}- -{}-parse-peg exemplos/busca.mpy}; sem
opção vale \texttt{-{}-parse-rd}. Os testes rodam com
\texttt{cabal test} ou \texttt{bash testes/rodar.sh}.

\subsection{Testes Realizados}

São 88 programas em \texttt{exemplos/}: os do repositório, os
fornecidos com o enunciado e os 15 da seção 9 da especificação. Os 27
com arquivo \texttt{.err} ao lado (entre eles os 15 de
\texttt{invalid/lexical} e \texttt{invalid/syntax}) são inválidos, e a
saída de erro de cada modo tem que ser exatamente o conteúdo do
arquivo. Nos outros 61, os três parsers têm que produzir a mesma AST,
posições incluídas; os de \texttt{invalid/semantic} só têm erros de
tipo e contam como válidos. Os 88 passam nos três modos.

\subsection{Dificuldades encontradas}
\label{sec:dificuldades}

\paragraph{Mensagens e posições do PEG.} Com a typed-peg 0.4 não havia
como apontar o erro, e a solução foi usar a 0.6, permitida pelo
enunciado. Mesmo assim, a primeira versão aceitava programas inválidos:
sob \texttt{gtR}, o \texttt{!.} do fim do programa tinha sucesso no meio
do arquivo, porque o primeiro caractere da linha seguinte era recusado
pela relação de coluna, e não por falta de entrada. Os testes com os
\texttt{.err} mostraram o problema, e o fim do arquivo passou a ser
testado sob \texttt{\_\textasciitilde}.

\paragraph{Ambiguidade da lambda.} O LALR teve 16 conflitos
shift/reduce depois do corpo de uma lambda, diante de um operador:
deslocar (o operador faz parte do corpo) ou reduzir. A primeira solução
foi um pseudo-token de menor precedência com \texttt{\%prec}, pois
\texttt{'lambda'} com baixa precedência não adianta: a regra herda a do
último terminal, o \texttt{:}. Mas os conflitos eram sintoma de uma
ambiguidade da própria GLC, com a lambda entre os átomos; com ela no
nível de \texttt{Expr}, nos três parsers, deixaram de existir.

\paragraph{Diferenças entre os parsers.} Os três liam \texttt{-2 ** 2}
como \texttt{-(2 ** 2)}, contra a especificação, e só a releitura dela
mostrou o erro. Comparar os três revelou problemas do LALR: aceitava
\texttt{1 if 2 if 3 else 4 else 5}, que os outros recusam, e recusava
\texttt{a if b else c if d else e} e o \texttt{pass} como corpo de
classe. Todos foram corrigidos.

\subsection{Limitações}

Em dois casos de erro fora dos testes, a mensagem do PEG difere da dos
outros parsers: uma palavra
reservada onde se espera um nome (\texttt{for if in xs:}) sai como
\texttt{unexpected 'if'}, e um parêntese não fechado na última linha é
apontado no \texttt{\textbackslash n} dessa linha, e não no
\texttt{NEWLINE} que o analisador léxico gera no fim do arquivo.

\section{Uso de IA generativa}

O registro dos prompts utilizados, das respostas obtidas e das
modificações feitas, pedido no item 4 da seção 8 do enunciado, está no
Apêndice~\ref{sec:ia}, fora do limite de seis páginas do relatório.

\section{Conclusão}

Os três parsers reconhecem a mesma linguagem e produzem a mesma AST. A
GLC estratificada levou quase direto ao descendente recursivo; no LALR
os conflitos revelaram a ambiguidade da lambda, e no PEG o trabalho foi
tratar léxico e indentação na própria gramática.

\begin{thebibliography}{9}
\setlength{\itemsep}{0pt}
\bibitem{appel1998}
Appel, A. W. (1998). \emph{Modern Compiler Implementation in ML}. Cambridge University Press.

\bibitem{ford2004}
Ford, B. (2004). Parsing Expression Grammars: A Recognition-Based Syntactic Foundation. In \emph{POPL '04}. ACM.

\bibitem{adams2013}
Adams, M. D. (2013). Principled Parsing for Indentation-Sensitive Languages: Revisiting Landin's Offside Rule. In \emph{POPL '13}. ACM.

\bibitem{typedpeg}
Ribeiro, R. (2026). \emph{typed-peg: Type-safe PEG parser combinators}, versão 0.6.0.0. Hackage.

\bibitem{megaparsec}
Karpov, M. \emph{megaparsec: Monadic parser combinators}, versão 9.8. Hackage.

\bibitem{happy}
Marlow, S.; Gill, A. \emph{Happy: The Parser Generator for Haskell}, versão 2.2.
\end{thebibliography}

\clearpage
\appendix
```latex
% ============================================================
% Seção: Uso de IA generativa
% Brenda Gabrielle Alves Nascimento, 24.1.4011
% BCC328 - Construção de Compiladores I - Trabalho Prático 1
% Inserir antes da conclusão.
% ============================================================

\section{Uso de IA generativa}
\label{sec:ia}

O desenvolvimento do trabalho contou com o apoio de uma ferramenta
de inteligência artificial generativa, o assistente Claude
(Anthropic), utilizada ao longo das diferentes etapas de implementação.
A ferramenta foi empregada como recurso de apoio à compreensão dos
requisitos, discussão de decisões de projeto, consulta sobre conceitos
e bibliotecas da linguagem Haskell, elaboração de trechos de código
e investigação de problemas identificados durante a compilação e os
testes.

A interação ocorreu de forma incremental. As solicitações eram
formuladas a partir dos requisitos do trabalho, das decisões de
implementação e dos problemas encontrados durante o desenvolvimento.
As respostas obtidas eram analisadas e comparadas com a especificação,
a documentação das bibliotecas e o material disponibilizado na
disciplina. Quando necessário, as sugestões eram corrigidas, adaptadas
ou descartadas antes de serem incorporadas ao projeto.

A ferramenta também auxiliou na elaboração de módulos e trechos de
código. Entretanto, sua utilização não substituiu a análise dos
requisitos nem a avaliação das soluções adotadas. As decisões de
projeto, a escolha entre alternativas, a integração dos componentes
e a validação do comportamento do compilador permaneceram sob
responsabilidade da autora.

Como parte do processo de verificação, foi adotado o ciclo de
desenvolvimento incremental: definir o problema ou a decisão de
projeto, consultar a ferramenta quando pertinente, analisar a solução
proposta, implementar ou adaptar o código, compilar com
\texttt{-Wall -Werror}, executar os testes e registrar as alterações
aprovadas no Git. Dessa forma, os resultados sugeridos pela IA eram
tratados como propostas a serem verificadas, e não como soluções
automaticamente corretas.

\subsection{Planejamento e organização do desenvolvimento}

\textbf{Contexto da interação.} No início do trabalho, foi necessário
organizar as etapas de implementação a partir do enunciado e da
especificação, considerando as dependências entre os componentes
do compilador.

\textbf{Apoio da IA.} Foi solicitado ao assistente um plano de
desenvolvimento incremental, contemplando desde a configuração do
ambiente até a elaboração do relatório e a validação final.

\textbf{Análise e decisões.} A ferramenta propôs um plano dividido
em dez etapas, com uma sequência baseada nas dependências entre
as fases do compilador. O plano foi utilizado como referência para
organizar o trabalho, mantendo-se a possibilidade de ajustar a ordem
e o escopo das atividades conforme os problemas encontrados.

\textbf{Validação.} Durante a configuração do ambiente, foi
identificada uma incompatibilidade entre a versão do GHC 8.8.4
disponibilizada inicialmente pela distribuição e a restrição
\texttt{base \^{}>=4.18.2.1} exigida pelo esqueleto da disciplina.
A instalação foi refeita utilizando GHCup, com o GHC 9.6.7,
versão compatível com o ambiente oficial da disciplina.

\subsection{Consulta e alinhamento com o material da disciplina}

\textbf{Contexto da interação.} Para manter a organização do projeto
compatível com os exemplos utilizados na disciplina, foi realizada
uma consulta ao repositório \texttt{bcc328-2026-2}, disponibilizado
pelo professor.

\textbf{Apoio da IA.} O assistente auxiliou na análise da estrutura
do repositório e na identificação de convenções presentes nos
exemplos, especialmente a organização dos diretórios
\texttt{Frontend/Lexer}, \texttt{Frontend/Parser},
\texttt{Frontend/Syntax} e \texttt{Frontend/Pretty}. Também foi
observada a organização do exemplo \texttt{Exp}, que mantém
implementações distintas de analisador léxico compartilhando
um único arquivo \texttt{Token.hs}.

\textbf{Análise e decisões.} A estrutura do projeto foi reorganizada
para seguir essas convenções. A classe \texttt{Pretty}, baseada em
\texttt{Text.PrettyPrint.HughesPJ}, foi adotada a partir do arquivo
\texttt{Utils/Pretty.hs} do repositório da disciplina.

Durante essa etapa, duas informações fornecidas pelo assistente
precisaram ser revistas. A primeira envolvia o atributo
\texttt{pkgs.haskell.compiler.ghc966}, utilizado no
\texttt{flake.nix} sugerido, que não estava disponível em
\emph{nixpkgs}; a consulta ao arquivo do professor indicou
\texttt{ghc96} como atributo adequado. A segunda dizia respeito
à possibilidade de o Alex manter estado entre linhas. A leitura
do arquivo \texttt{LineLexer.x} mostrou que o wrapper
\texttt{monadUserState} permite manter estado durante a análise.

A decisão de utilizar um \emph{scanner} \emph{ad-hoc} foi mantida,
mas por outro motivo: uma ação do Alex devolve um único
\emph{token}, enquanto a emissão de múltiplos \texttt{DEDENT}
consecutivos exigiria uma fila de \emph{tokens} pendentes no estado
do analisador.

\textbf{Validação.} As convenções adotadas foram comparadas com
os arquivos do repositório da disciplina, e as informações
divergentes foram corrigidas antes de serem utilizadas no projeto.

\subsection{Apoio na implementação do analisador léxico e da camada de indentação}

\textbf{Contexto da interação.} A implementação do analisador léxico
exigia definir as estruturas de posição, lexemas e tokens, além
de tratar as regras específicas de indentação da linguagem.

\textbf{Apoio da IA.} O assistente auxiliou na elaboração dos tipos
\texttt{Pos}, \texttt{Lexeme} e \texttt{Token}, do módulo
\texttt{MiniPythonLexer.hs} e da camada de indentação
\texttt{Indent.hs}. A implementação da indentação foi discutida
em etapas, inicialmente com a pilha de níveis e, depois, com
as exceções descritas na Seção 4.5 da especificação.

\textbf{Análise e decisões.} As palavras \texttt{True} e
\texttt{False} foram retiradas da lista de palavras reservadas
e passaram a ser reconhecidas como \texttt{TBoolLit}, evitando
duas representações para o mesmo valor. Posteriormente, foi
necessário derivar \texttt{Ord} em \texttt{Lexeme} e
\texttt{Token}, requisito da instância \texttt{Stream} do
Megaparsec.

A implementação da camada de indentação em etapas permitiu
identificar um problema na primeira versão: um programa com
linha em branco no meio de um bloco produzia sequências
indevidas de \texttt{DEDENT} e \texttt{INDENT}. Esse resultado
indicou a necessidade de considerar o tratamento previsto no
item 4 da Seção 4.5 antes de finalizar a implementação.

\textbf{Validação.} Os módulos foram compilados e testados com
programas que exercitavam as regras de indentação, incluindo
casos com linhas em branco e encerramento de blocos. Os
comportamentos divergentes foram utilizados para orientar
as correções.

\subsection{Modelagem da AST e representação textual}

\textbf{Contexto da interação.} A construção da árvore sintática
abstrata (AST) deveria seguir a Seção 3 do enunciado, armazenando
a posição de cada nó. Também era necessário implementar uma
representação textual própria, sem depender de instâncias
\texttt{Show} geradas automaticamente.

\textbf{Apoio da IA.} O assistente auxiliou na elaboração dos
módulos \texttt{MiniPythonSyntax.hs} e
\texttt{MiniPythonPretty.hs}, incluindo a organização dos tipos
da AST e a impressão de expressões considerando os níveis
de precedência.

\textbf{Análise e decisões.} A primeira versão apresentou
problemas de compilação relacionados a instâncias órfãs.
Como \texttt{-Wall} inclui o aviso \texttt{-Worphans}, foi
necessário tratar explicitamente essa condição. A solução
adotada utilizou \texttt{\{-\# OPTIONS\_GHC -Wno-orphans \#-\}}
acompanhada de comentário justificando a separação entre
a classe, os tipos e as instâncias de \texttt{Pretty}. Essa
abordagem também aparece no arquivo \texttt{LineLexer.x}
do professor.

Outros ajustes envolveram o acesso a um registro, em que
\texttt{case clsSuper of} foi corrigido para
\texttt{case clsSuper cd of}, e a composição da saída textual:
o uso de \texttt{punctuate (text "")} com \texttt{vcat}
não produzia a linha em branco esperada, sendo substituído
por \texttt{\$+\$}.

\textbf{Validação.} A AST e o impressor foram verificados por
meio da compilação e da comparação da representação gerada
com a estrutura esperada para os programas de teste.

\subsection{Discussão da gramática e implementação do analisador descendente recursivo}

\textbf{Contexto da interação.} Foi necessário definir uma
gramática livre de contexto (GLC) não ambígua e implementar
um analisador descendente recursivo organizado em módulos.

\textbf{Apoio da IA.} O assistente auxiliou na discussão das
produções da gramática e na elaboração de uma primeira versão
de \texttt{docs/gramatica.md}, além de apoiar a organização
do analisador nos módulos \texttt{Core.hs}, \texttt{Expr.hs}
e \texttt{RecursiveDescent.hs}.

\textbf{Análise e decisões.} A primeira versão da gramática
precisou ser revista porque a produção \texttt{Alvo} permitia
chamadas de função como alvo de atribuição, deixando essa
restrição apenas em uma explicação textual. A produção foi
reescrita para expressar a restrição diretamente na gramática.

Na implementação do analisador, os quatro programas de teste
falharam inicialmente em expressões como \texttt{int(input())}.
A investigação mostrou que \texttt{int}, \texttt{float},
\texttt{str} e \texttt{bool} são palavras reservadas de tipo,
mas também funcionam como conversões primitivas na Seção 8
da especificação. O não-terminal \texttt{Atomo} não tratava
esses identificadores em posição de expressão, sendo necessário
ajustar a análise sintática.

Também foi revista a convenção de posição dos nós. A primeira
versão utilizava a posição do operador em \texttt{EBin}; por
isso, um erro em \texttt{1 + 2 = 3} apontava para o símbolo
\texttt{+}. Foi adotada a convenção de que a posição de um nó
corresponde ao início da construção no texto-fonte, o que
também favorece a localização de erros em etapas posteriores,
como a verificação de tipos.

\textbf{Validação.} A gramática e o analisador foram avaliados
com os programas de teste da especificação, incluindo casos
válidos e inválidos. As falhas observadas orientaram as
alterações nas produções e nas regras de análise.

\subsection{Comparação de alternativas e implementação do analisador PEG}

\textbf{Contexto da interação.} A implementação do analisador PEG
exigia a utilização da biblioteca \texttt{typed-peg}. Diferentemente
do analisador descendente recursivo, que utiliza uma lista de tokens
produzida pelo analisador léxico, o analisador PEG realiza a análise
diretamente sobre o texto-fonte.

\textbf{Apoio da IA.} O assistente auxiliou na compreensão da
estrutura da biblioteca, na elaboração das regras da gramática PEG
e na implementação do módulo \texttt{PEG.hs}, utilizando o
quasi-quoter \texttt{pegGrammar}. Também contribuiu na investigação
de problemas relacionados ao reconhecimento das construções da
linguagem e à localização dos erros no código-fonte.

\textbf{Análise e decisões.} Durante o desenvolvimento, foi
avaliado o uso da implementação PEG disponibilizada na disciplina
(\texttt{workspace/src/PEG/}). Após a análise de sua estrutura
e dos recursos utilizados, optou-se pela biblioteca
\texttt{typed-peg}, alternativa permitida pelo enunciado.

A implementação exigiu ajustes na gramática para contemplar
o reconhecimento das construções da linguagem diretamente no
texto-fonte, incluindo o tratamento da indentação. Também foi
necessário considerar as informações de posição disponibilizadas
pela biblioteca para produzir mensagens de erro adequadas.

A utilização do \texttt{typed-peg} foi mantida na versão final
do projeto. É importante destacar que o Megaparsec foi utilizado
no analisador descendente recursivo, e não no PEG. Essa separação
permitiu manter as duas implementações sintáticas independentes,
conforme a organização proposta para o trabalho.

\textbf{Validação.} O analisador PEG foi compilado e executado
com os programas de teste da especificação. Seus resultados
foram comparados com os produzidos pelos outros analisadores,
verificando o reconhecimento das construções da linguagem e
a consistência dos resultados.

\subsection{Implementação e depuração do analisador LALR}

\textbf{Contexto da interação.} A implementação do analisador
LALR com Happy deveria declarar a precedência por meio de
diretivas, sem estratificar as produções de expressão,
conforme a Seção 5 do enunciado.

\textbf{Apoio da IA.} O assistente auxiliou na elaboração
inicial do arquivo \texttt{Grammar.y} e na investigação
dos conflitos reportados pelo gerador Happy.

\textbf{Análise e decisões.} A primeira versão apresentou
16 conflitos \emph{shift/reduce}. A investigação utilizou
o arquivo de informações gerado pelo comando
\texttt{happy --info}, que indicou que os conflitos estavam
concentrados em um único estado e relacionados à produção
do corpo da \emph{lambda}.

A primeira tentativa de resolução sugerida, por meio da
declaração \texttt{\%right 'lambda'}, não produziu efeito.
A análise mostrou que, no Happy, a precedência de uma
produção é herdada de seu último terminal, que naquela
regra era o símbolo \texttt{:}. A solução adotada foi
introduzir o pseudo-\emph{token} \texttt{LAMBDA}, com a
menor precedência da tabela, e utilizar
\texttt{\%prec LAMBDA} para sobrescrever a precedência
da produção. Com essa alteração, os conflitos foram
eliminados.

Também foram identificadas dependências ausentes na
configuração do projeto: \texttt{array}, exigida pelo
código gerado, e \texttt{build-tool-depends: happy:happy}
no arquivo \texttt{.cabal}.

\textbf{Validação.} A saída de diagnóstico do Happy foi
utilizada para localizar os conflitos, e a geração do
analisador foi repetida após as alterações, até que
não fossem mais reportados conflitos.

\subsection{Apoio na elaboração e validação dos testes}

\textbf{Contexto da interação.} Era necessário construir
programas de teste com base nos exemplos da Seção 9 da
especificação e verificar se os três analisadores
apresentavam resultados equivalentes.

\textbf{Apoio da IA.} O assistente auxiliou na elaboração
de nove programas válidos, sete programas inválidos e
dos \emph{scripts} \texttt{testes/rodar.sh} e
\texttt{testes/gerar-saidas.sh}, utilizados na execução
e comparação dos analisadores.

\textbf{Análise e decisões.} A primeira execução do
comparativo apresentou falsos positivos: como a compilação
estava quebrada, os dois lados da comparação retornavam
saídas vazias, que eram interpretadas pelo \texttt{diff}
como iguais. O \emph{script} foi então modificado para
armazenar as saídas em arquivos e verificar se elas
não estavam vazias antes de realizar a comparação.

\textbf{Validação.} Os testes foram executados novamente
após a correção do \emph{script}. A verificação passou
a considerar não apenas a igualdade entre as saídas,
mas também se os programas haviam sido executados
corretamente e produzido resultados.

\subsection{Limitações e avaliação das respostas}

Durante o desenvolvimento, algumas respostas fornecidas
pela ferramenta apresentaram inconsistências ou não
consideraram completamente as particularidades do ambiente
e da especificação da disciplina. Entre os casos
identificados estavam a indicação de um atributo
inexistente na configuração do Nix, uma interpretação
incorreta sobre o funcionamento do Alex e uma tentativa
inicial de resolver os conflitos de precedência do Happy
que não produziu o efeito esperado.

Essas situações evidenciaram a importância de avaliar
criticamente as sugestões recebidas. Para isso, foram
consultados o material disponibilizado pelo professor,
a documentação das bibliotecas e os resultados obtidos
com a compilação e os testes. As sugestões que não
atendiam aos requisitos foram corrigidas, adaptadas
ou descartadas.

A experiência mostrou que a IA foi útil como apoio
à investigação e à implementação, especialmente para
explorar alternativas e acelerar a identificação de
possíveis causas de erros. Entretanto, a correção das
soluções dependia de sua verificação no contexto real
do projeto.

\subsection{Síntese}

A IA generativa foi utilizada como assistente técnico
durante diferentes etapas do trabalho. Sua contribuição
envolveu a organização do desenvolvimento, a consulta
sobre conceitos e bibliotecas, a discussão de alternativas
de projeto, a elaboração de trechos de código e o apoio
na investigação de falhas.

A participação da ferramenta foi especialmente útil
em tarefas que exigiram integração entre componentes,
como o tratamento de indentação, a construção dos
analisadores sintáticos e a resolução de conflitos
de precedência no Happy. Ao mesmo tempo, as interações
também permitiram identificar limitações nas respostas
e a necessidade de confrontá-las com a especificação
e com o material da disciplina.

O desenvolvimento foi conduzido de forma incremental,
com compilação utilizando \texttt{-Wall -Werror},
execução de testes e registro das alterações no Git.
As sugestões da ferramenta não foram incorporadas
automaticamente: foram analisadas, adaptadas e
validadas de acordo com os requisitos do trabalho.
Assim, a IA atuou como recurso de apoio à implementação
e à resolução de problemas, enquanto as decisões de
projeto, a análise crítica e a validação final do
compilador permaneceram sob responsabilidade da autora.
```

\end{document}

